\section{DISCUSSION}
This paper systematically reviews the Python type analysis ecosystem and
compares how static, dynamic, and AI-based approaches address specific
challenges while highlighting their limitations.
Through this, the impact of Python’s dynamic features on the design of analysis
techniques becomes evident.

As observed in RQ1, achieving program safety with a single static type system
in Python is challenging.
Consequently, static type checking, runtime enforcement, and hybrid approaches
coexist, forming a unique technical landscape.
This layered strategy represents a practical compromise, balancing developer
productivity and dynamic flexibility while ensuring a certain level of
reliability and error detection.

The structural challenges highlighted in RQ2 introduce problems that cannot be
fully addressed by traditional language theory.
Specifically, structural typing, flow-sensitive type changes, heterogeneous
containers, and dynamic code generation via reflection or eval fundamentally
challenge assumptions of conventional static analysis (e.g., nominal hierarchy,
stable object model, closed world assumption).
These issues demand reconsideration of the analysis model itself.
For instance, flow-sensitive analyses provide high precision but incur high
computational costs, whereas flow-insensitive approaches scale better but with
limited precision.
This underscores the constant trade-offs Python analyzers must navigate among
precision, scalability, and soundness.

A similar trend emerges in the tool comparison of RQ3. Static analyzers remain
essential for large-scale codebases, yet they cannot fully model Python’s
dynamic features.
Dynamic and hybrid tools improve accuracy using runtime information but cannot
guarantee fundamental completeness due to coverage limitations.
Research-oriented tools leveraging constraint solving or symbolic execution
offer high precision but are limited in practical adoption due to computational
cost and challenges in handling dynamic code.
Recently, LLM-based analyses provide high inference coverage even in
low-annotation settings but cannot ensure soundness due to non-determinism and
hallucinations.

Overall, the central challenge in Python type analysis is the discontinuity
between the language’s dynamic, open-ended nature and the demands of static
safety.
Each approach provides strategies to bridge this gap, yet no single method
offers a complete solution.
Hybrid or feedback-driven approaches that integrate static, dynamic, and
AI-based techniques are increasingly recognized as a promising direction for
future development.
